% ============================================================
% 40_colors_structure.tex — Farbpaletten & Kapitel-Orchestrierung
% ============================================================

% Index-Palette (18 Slots, kuratiert fuer 3-Fenster-Kontrast)
% ------------------------------------------------------------
% Reihenfolge: ETH-Signaturfarben zuerst, danach global distanzierte
% Zusatzfarben. Benachbarte Kapitel (Abstand 1 UND 2) bleiben deutlich
% unterscheidbar; globale Fast-Dopplungen werden vermieden, soweit es mit
% 18 kontraststarken Slots und fixierten ETH-Seeds sinnvoll moeglich ist.
% Slot 0 ist reserviert fuer den Frontchapter (\StartFrontChapter) -- ETH-Blau.
% Niemals ChapterColor* in Kapiteldateien hart-coden; ueber \chaptercolor /
% \chaptercolorlight / \chaptercolortext referenzieren.
\definecolor{ChapterColor0}{HTML}{215CAF}       % ETH Blue
\definecolor{ChapterColor0Light}{HTML}{CEDBED}
\definecolor{ChapterColor0Text}{HTML}{FFFFFF}
\definecolor{ChapterColor1}{HTML}{6E3B6E}       % ETH Purple
\definecolor{ChapterColor1Light}{HTML}{DFD4DF}
\definecolor{ChapterColor1Text}{HTML}{FFFFFF}
\definecolor{ChapterColor2}{HTML}{5A7F2B}       % ETH Green
\definecolor{ChapterColor2Light}{HTML}{DBE3D0}
\definecolor{ChapterColor2Text}{HTML}{FFFFFF}
\definecolor{ChapterColor3}{HTML}{8C6239}       % ETH Bronze
\definecolor{ChapterColor3Light}{HTML}{E6DCD3}
\definecolor{ChapterColor3Text}{HTML}{FFFFFF}
\definecolor{ChapterColor4}{HTML}{007894}       % ETH Petrol
\definecolor{ChapterColor4Light}{HTML}{C7E1E7}
\definecolor{ChapterColor4Text}{HTML}{FFFFFF}
\definecolor{ChapterColor5}{HTML}{B7352D}       % ETH Red
\definecolor{ChapterColor5Light}{HTML}{EFD3D1}
\definecolor{ChapterColor5Text}{HTML}{FFFFFF}
\definecolor{ChapterColor6}{HTML}{23523B}       % Dark Spruce
\definecolor{ChapterColor6Light}{HTML}{CFD9D4}
\definecolor{ChapterColor6Text}{HTML}{FFFFFF}
\definecolor{ChapterColor7}{HTML}{8F2459}       % Raspberry
\definecolor{ChapterColor7Light}{HTML}{E6CFDA}
\definecolor{ChapterColor7Text}{HTML}{FFFFFF}
\definecolor{ChapterColor8}{HTML}{17205E}       % Midnight Blue
\definecolor{ChapterColor8Light}{HTML}{CCCEDC}
\definecolor{ChapterColor8Text}{HTML}{FFFFFF}
\definecolor{ChapterColor9}{HTML}{2D8659}       % Emerald
\definecolor{ChapterColor9Light}{HTML}{D1E4DA}
\definecolor{ChapterColor9Text}{HTML}{FFFFFF}
\definecolor{ChapterColor10}{HTML}{523523}      % Walnut
\definecolor{ChapterColor10Light}{HTML}{D9D3CF}
\definecolor{ChapterColor10Text}{HTML}{FFFFFF}
\definecolor{ChapterColor11}{HTML}{8F2481}      % Orchid
\definecolor{ChapterColor11Light}{HTML}{E6CFE3}
\definecolor{ChapterColor11Text}{HTML}{FFFFFF}
\definecolor{ChapterColor12}{HTML}{1D581D}      % Deep Green
\definecolor{ChapterColor12Light}{HTML}{CDDACD}
\definecolor{ChapterColor12Text}{HTML}{FFFFFF}
\definecolor{ChapterColor13}{HTML}{5E1720}      % Oxblood
\definecolor{ChapterColor13Light}{HTML}{DCCCCE}
\definecolor{ChapterColor13Text}{HTML}{FFFFFF}
\definecolor{ChapterColor14}{HTML}{4C248F}      % Violet
\definecolor{ChapterColor14Light}{HTML}{D8CFE6}
\definecolor{ChapterColor14Text}{HTML}{FFFFFF}
\definecolor{ChapterColor15}{HTML}{233552}      % Slate Navy
\definecolor{ChapterColor15Light}{HTML}{CFD3D9}
\definecolor{ChapterColor15Text}{HTML}{FFFFFF}
\definecolor{ChapterColor16}{HTML}{207E73}      % Sea Teal
\definecolor{ChapterColor16Light}{HTML}{CEE3E0}
\definecolor{ChapterColor16Text}{HTML}{FFFFFF}
\definecolor{ChapterColor17}{HTML}{58581D}      % Olive Brown
\definecolor{ChapterColor17Light}{HTML}{DADACD}
\definecolor{ChapterColor17Text}{HTML}{FFFFFF}

% ------------------------------------------------------------
% Warn-/Semantik-Farben
% ------------------------------------------------------------
\definecolor{ETHRot}{HTML}{CC0000}
\colorlet{WarnboxBack}{ETHRot!15}
\colorlet{DangerBackColor}{ETHRot!85!black}
\colorlet{DangerTextColor}{white}

% Semantisches Formel-Highlighting (kapitelunabhängig, siehe \ZSFmhlA/B/C/D).
\definecolor{MathHLPrimary}{HTML}{00796B}    % Quelle / gegeben / erster Strang
\definecolor{MathHLSecondary}{HTML}{C2410C}  % Gegenstueck / abgeleitet / zweiter Strang
\definecolor{MathHLTertiary}{HTML}{A50D68}   % dritter paralleler Strang (sparsam)
\definecolor{MathHLTarget}{HTML}{1D4ED8}     % Ziel / Endform / Resultat

% Farben für textVorBox / textNachBox (Text-Box-Bindungen)
\definecolor{TextVorBoxColor}{HTML}{0A0A1E}  % Sehr dunkles Blau (einleitend)
\definecolor{TextNachBoxColor}{HTML}{0A1A0A}  % Sehr dunkles Grün (ergänzend)

% ------------------------------------------------------------
% Dokumenttitel-Masthead: bewusst unabhängig von der Kapitelrotation.
\colorlet{ZSFDocumentTitleBarColor}{ChapterColor0}
\colorlet{ZSFDocumentTitleTextColor}{ChapterColor0Text}
\colorlet{ZSFDocumentTitleBylineColor}{ChapterColor0Text!88}

% Schlusswörter in eingebetteten Kapiteln: ETH-Palette für den Inhalt,
% aber dieselbe klare Leistenfarbe wie der Dokumenttitel.
\newcommand{\ZSFStartClosingSubsection}{%
  \ZSFSetChapterPaletteByIndex{0}%
  \renewcommand{\SubsectionBarColor}{ZSFDocumentTitleBarColor}%
  \renewcommand{\SubsectionBarTextColor}{ZSFDocumentTitleTextColor}%
}

% Aktive Kapitelfarben (werden pro Kapitel umgeschaltet)
% ------------------------------------------------------------
\newcommand{\chaptercolor}{ChapterColor1}
\newcommand{\chaptercolorlight}{ChapterColor1Light}
\newcommand{\chaptercolortext}{ChapterColor1Text}
\newcommand{\ChapterBarColor}{\chaptercolor}
\newcommand{\ChapterBarTextColor}{\chaptercolortext}
\newcommand{\SubsectionBarColor}{\chaptercolor!75}
\newcommand{\SubsectionBarTextColor}{\chaptercolortext}
\newcommand{\SubsubsectionBarColor}{\chaptercolorlight!78!white}
\newcommand{\SubsubsectionTitleColor}{\chaptercolor!72!black}
\newcommand{\FormulaboxBackColor}{\chaptercolorlight!40}
\newcommand{\ZSFtableFrameColor}{\chaptercolor}         % Duenner Aussenrahmen der Tabellen

% Friert die Tabellen-Rahmenfarbe auf die AKTUELL aktive Kapitelfarbe ein.
% Gedacht fuer Box-Gruppen, die als ein Block gelesen werden, deren Boxen
% aber intern verschiedene Kapitelpaletten fuehren (Symbol-Tabellen in
% ch00): innen bleibt jede Box kapitelfarbig, der Rahmen bleibt einheitlich.
% Wirkt nur bis zum Ende der umgebenden TeX-Gruppe.
\newcommand{\ZSFPinTableFrameColor}{\edef\ZSFtableFrameColor{\chaptercolor}}
\newcommand{\BoxTitleColor}{\chaptercolorlight}
\newcommand{\BoxTitleTextColor}{TextVorBoxColor}        % Titeltext-Farbe (dunkel auf Pastell)
\newcommand{\BoxBackColor}{\chaptercolorlight!8}        % Sehr sanfter Box-Hintergrund
\newcommand{\BoxFrameColor}{\chaptercolorlight!25}      % Rahmen in der Farbe des vorherigen Hintergrunds

% Aussage-, Ziel- und Grid-Boxen.
\newcommand{\StatementBarColor}{\chaptercolor}
\newcommand{\StatementTitleColor}{\chaptercolor}
\newcommand{\GoalConditionBackColor}{white}
\newcommand{\GoalConditionFrameColor}{\chaptercolorlight!22}
\newcommand{\GoalConditionBorderColor}{black}
\newcommand{\GoalTargetBackColor}{\chaptercolorlight!22}
\newcommand{\GoalTargetFrameColor}{\chaptercolor}
\newcommand{\CompactGridFrameColor}{black!20}

% Index-basierte Palette (automatisch via \StartChapter, Loop mit Wrap-Warnung)
\newcount\ZSFchapterPaletteIndex
\newcount\ZSFchapterPaletteLookupIndex
\newcommand{\ZSFchapterPaletteSize}{18}       % Anzahl verfügbarer Indexfarben (0..17)
\newcommand{\ZSFSetChapterPaletteByIndex}[1]{%
  \renewcommand{\chaptercolor}{ChapterColor#1}%
  \renewcommand{\chaptercolorlight}{ChapterColor#1Light}%
  \renewcommand{\chaptercolortext}{ChapterColor#1Text}%
}
\newcommand{\ZSFSetPaletteColorNameFromNumber}[2]{%
  \ZSFchapterPaletteLookupIndex=#1\relax
  \loop
  \ifnum\ZSFchapterPaletteLookupIndex>\numexpr\ZSFchapterPaletteSize-1\relax
    \advance\ZSFchapterPaletteLookupIndex by -\ZSFchapterPaletteSize
  \repeat
  \edef#2{ChapterColor\the\ZSFchapterPaletteLookupIndex}%
}
\newcommand{\SetChapterPaletteByNumber}{%
  \ZSFchapterPaletteIndex=\value{section}%
  \loop
  \ifnum\ZSFchapterPaletteIndex>\numexpr\ZSFchapterPaletteSize-1\relax
    \advance\ZSFchapterPaletteIndex by -\ZSFchapterPaletteSize
  \repeat
  \ifnum\value{section}>\numexpr\ZSFchapterPaletteSize-1\relax
    \PackageWarning{ZSF}{Chapter palette wrapped at section \thesection. Increase palette if unique colors are required.}%
  \fi
  \ZSFSetChapterPaletteByIndex{\the\ZSFchapterPaletteIndex}%
}

% ------------------------------------------------------------
% Chapter-Start-Orchestrierung
% ------------------------------------------------------------
\newif\ifzsfAfterChapterBar
\zsfAfterChapterBarfalse
\newcommand{\ZSFConsumeAfterChapterBar}{%
  \global\zsfAfterChapterBarfalse
}
\newcommand{\ZSFMarkAfterChapterBar}{%
  \global\zsfAfterChapterBartrue
  \AddToHookNext{para/begin}{\ZSFConsumeAfterChapterBar}%
}
% Kollabiert den before-skip der SubsectionBar auf \ZSFsubsectionAfterChapterGap,
% wenn sie direkt nach einer ChapterBar steht. Reset nach Verwendung.
\newcommand{\ZSFCollapseAfterChapterBar}{%
  \ifzsfAfterChapterBar
    \vspace*{\dimexpr\ZSFsubsectionAfterChapterGap-\ZSFsubsectionBarBeforeSkip\relax}%
    \ZSFConsumeAfterChapterBar
  \fi
}

% ------------------------------------------------------------
% SubsectionBar-Orchestrierung (analog zur ChapterBar oben):
% Eine Box direkt nach einer SubsectionBar kollabiert ihren before-skip auf
% \ZSFsubsectionAfterGap, damit sie an der Bar klebt. Das Flag wird beim
% nächsten Absatzbeginn automatisch gelöscht (wie bei der ChapterBar).
% ------------------------------------------------------------
\newif\ifzsfAfterSubsectionBar
\zsfAfterSubsectionBarfalse
\newcommand{\ZSFConsumeAfterSubsectionBar}{%
  \global\zsfAfterSubsectionBarfalse
}
\newcommand{\ZSFMarkAfterSubsectionBar}{%
  \global\zsfAfterSubsectionBartrue
  \AddToHookNext{para/begin}{\ZSFConsumeAfterSubsectionBar}%
}
% Box-before-skip: kollabiert nach einer SubsectionBar, sonst Standard-Luft
% der Titelboxen (\ZSFtitleboxBeforeSkip).
\newcommand{\ZSFBoxBeforeSkip}{%
  \ifzsfAfterSubsectionBar\ZSFsubsectionAfterGap\else\ZSFtitleboxBeforeSkip\fi
}

% Jedes Kapitel bekommt automatisch ein Seiten-Label `zsfchap:<nr>`, damit
% Register (z. B. die Symbol-Tabellen in ch00) die Startseite eines Kapitels
% nennen koennen, ohne dass in jeder Kapiteldatei ein Label gepflegt werden
% muss. Der Zugriff laeuft ueber \ZSFChapterPage.
\newcommand{\ZSFchapPageLabel}{}
\newcommand{\ZSFMarkChapterPage}{%
  \edef\ZSFchapPageLabel{zsfchap:\the\value{section}}%
  \expandafter\label\expandafter{\ZSFchapPageLabel}%
}
% Bewusst \pageref* (ohne Hyperlink): im Titelbalken würde die globale
% Linkfarbe als Fremdkörper lesen; so erbt die Zahl die Titeltext-Farbe.
\newcommand{\ZSFChapterPage}[1]{\pageref*{zsfchap:#1}}

\makeatletter
\newcommand{\ZSFPostChapterSpacing}{\vspace*{\ZSFpostChapterNudge}}
\newcommand{\StartChapter}[2][]{%
  \ifnum\value{zsfFirstChap}=1\setcounter{zsfFirstChap}{0}\fi
  \refstepcounter{section}
  \ZSFMarkChapterPage
  \SetChapterPaletteByNumber
  \setcounter{subsection}{0}
  \if\relax\detokenize{#1}\relax\else\label{#1}\fi
  \ChapterBar{\thesection. #2}
  \ZSFMarkAfterChapterBar
  \ZSFPostChapterSpacing
}
\makeatother

% Variante: Kapitel explizit auf neuer Spalte beginnen (wie Analysis)
\newcommand{\StartChapterOnNewColumn}[2][]{%
  \columnbreak
  \StartChapter[#1]{#2}%
}

\newcommand{\StartFrontChapter}[1]{%
  \ifnum\value{zsfFirstChap}=1\setcounter{zsfFirstChap}{0}\else\columnbreak\fi
  \renewcommand{\chaptercolor}{ChapterColor0}%
  \renewcommand{\chaptercolorlight}{ChapterColor0Light}%
  \renewcommand{\chaptercolortext}{ChapterColor0Text}%
  \setcounter{subsection}{0}%
  \ChapterBar{#1}%
  \ZSFMarkAfterChapterBar
  \ZSFPostChapterSpacing
}
